%!TEX root =  CV_Master.tex
%**
% =====================================================
% Part I — Compact summary
% =====================================================

% ----------------------------------------------------------------------
\cvsection{Main Academic Appointments}

\cvsubsection{Current}
\cventry{2024–}{Director, Interdisciplinary Lab (ILAS)}{International School for Advanced Studies (SISSA), Trieste, Italy}
\cventry{2023–}{Full Professor of Theoretical Physics (Professore di I fascia)}{Physics Department, SISSA, Trieste, Italy}

\cvsubsection{Previous academic positions}
\cventry{2021–25}{Head, Theoretical and Scientific Data Science group}{Physics Department, SISSA, Trieste, Italy}
\cventry{2019–22}{Full Professor of Astrostatistics (on leave 2020–22)}{Physics Department, Imperial College London, UK}
\cventry{2015–20}{Director, Centre for Languages, Culture and Communication}{Education Office, Imperial College London, UK}
\cventry{2008–19}{Lecturer (2008–12), Senior Lecturer (2012–16), Reader (2016–19)}{Physics Department, Imperial College London, UK}
\cventry{2006–08}{Junior Fellow}{St Anne's College, University of Oxford, UK}
\cventry{2005–08}{Sir Norman Lockyer Research Fellow of the Royal Astronomical Society}{Astrophysics Department, University of Oxford, UK}
\cventry{2004–05}{Postdoctoral Researcher}{Department of Theoretical Physics, University of Geneva, Switzerland}

\cvsubsection{Visiting Academic Positions}
\cventry{2023–}{Visiting Professor of Astrostatistics}{Imperial College London, UK}
\cventry{2023–}{Visiting Faculty, Data-Driven Management}{MIB Trieste School of Management, Italy}
\cventry{2019–22}{Visiting Professor of Cosmology}{Gresham College, London, UK}
\cventry{2011–13}{Next Einstein Visiting Scholar}{African Institute for Mathematical Sciences (AIMS), Cape Town, South Africa}

% ----------------------------------------------------------------------
\cvsection{Education and qualifications}
\cventry{2022, 2020}{Abilitazione Scientifica Nazionale a professore di I fascia}{Settori FIS 02/A2 (valid until 2032) and FIS 02/C1 (valid until 2029)}
\cventry{2011}{Certificate of advanced study in learning and teaching}{Centre for Higher Education Research and Scholarship, Imperial College London}
\cventry{2004}{PhD in Theoretical Physics}{University of Geneva, Switzerland. Supervisor: Ruth Durrer}
\cventry{2001}{MSc(Hons) in Physics}{Swiss Federal Institute of Technology (ETHZ), Zurich, Switzerland}


% ----------------------------------------------------------------------
\cvsection{Selected Prizes and Awards}

\textbf{Charmian Brinson Honorary Lecture}, Imperial College London (2022); \textbf{Annie Maunder Medal for Outreach}, Royal Astronomical Society (2020) --- the highest RAS recognition for public engagement;\textbf{President's Award for Leadership in Public Engagement}, Imperial College London (2018); \textbf{Chaire Georges Lema\^{\i}tre} (2018), Universit\'e Catholique de Louvain; \textbf{Royal Statistical Society Significance Lecture} (2017); \textbf{Outstanding Performance Award}, Imperial College (2017); \textbf{President's Award for Excellence in Teaching}, Imperial College London (2016); \textbf{Foreign Policy 100 Global Thinkers} (2014); \textbf{Michelson Postdoctoral Lectureship Prize}, Case Western Reserve University (2008); \textbf{Lord Kelvin Award Lecture}, British Association for the Advancement of Science (2007).

% ----------------------------------------------------------------------
\cvsection{Scientific Publications Record}

\cvsmallheading{Output.}
\refcvcount{nAllRefPapers} papers in peer-reviewed physics journals or machine learning conferences (12 of which as first author, including a single-author review article, 3 in \PRL, 1 in {\em Nature Astronomy}, 2 in NeurIPS, one as spotlight); \refcvcount{nProceedings} conference proceedings; \refcvcount{nWhitePapers} white papers and editorial works; \refcvcount{nBookChapters} book chapters.
The single-author review {\itshape `Bayes in the sky'} (Contemporary Physics, 2008) has 1500+ citations.

\cvsmallheading{Citations (October 2026).}
15{,}000+ citations, h-index 54 (\href{https://inspirehep.net/literature?sort=mostrecent&size=25&page=1&q=a%20R.Trotta.1&ui-citation-summary=true}{INSPIRE}/\href{https://ui.adsabs.harvard.edu/search/q=orcid%3A0000-0002-3415-0707&sort=date%20desc%2C%20bibcode%20desc&p_=0}{NASA ADS}, higher of each per paper); 13{,}000+ citations, h-index 57 (\href{https://scholar.google.com/citations?user=V5FvaSYAAAAJ}{Google Scholar}). Citation thresholds: 21 papers $>\!100$ citations; 2 papers $>\!500$; 3 papers $>\!1000$. {\itshape Full publications list given in Part II.}

\cvsmallheading{Major community and white papers.}
{\itshape `Opportunities in AI/ML for the Rubin LSST Dark Energy Science Collaboration'} White Paper (2026);
{\itshape Nature Astronomy} opinion piece (2025): `The indiscriminate adoption of AI threatens the foundations of academia'; co-lead of the Astrostatistics section, EuCAPT White Paper (2021); Astrostatistics section lead, Euclid Theory Working Group {\itshape `Cosmology and Fundamental Physics with the Euclid Satellite'}, Living Reviews in Relativity (2013, 2018).


% ----------------------------------------------------------------------
\cvsection{Competitive Grants Record}

Full list and per-section breakdown in Part II. Figures below are the running total across all competitive grants held as PI or co-I (excludes the one declined award).

\grantgrandtotal


% ----------------------------------------------------------------------
\cvsection{Main Leadership Roles and Responsibilities}

\cvsmallheading{SISSA, Trieste.}
Director, Interdisciplinary Lab (ILAS, 2024–); Head of Theoretical and Scientific Data Science group (2021–25): co-funder of the group and responsible for setting up the associated PhD programme; Member, Library committee; Member, Physics Area Steering Board (2021–25); Knowledge Transfer Committee representative (2021–); Physics Area coordinator for the 2014–19 national assessment exercise (VQR, 2021).
\cvsmallheading{Imperial College London.}
Director, Centre for Languages, Culture and Communication (CLCC, 2015–20): responsible for {\em Imperial Horizons}, encompassing 180+ modules reaching 4000+ undergraduates; led a full curriculum review for Imperial Horizons (2018) and one of the architects of the wider {\em I-Explore} co-curricular offering for all Imperial undergraduates (2019); member of Imperial College Senate (2016–20); Education Strategy and Operations Group (2020); Education Office COVID-19 Response Group (2020); Chair, Curriculum Review Reference Panel (2018–19); Chair of CLCC Departmental Teaching Committee, CLCC Management Team, Imperial Horizons Team, and President's Award review panel (all 2015–20); co-chair, CLCC/CHERS Education Committee (2016–20).

\cvsmallheading{Gresham College, London.}
Member of Academic Board (2019–22).

\cvsmallheading{International.}
Member, EuCAIF Steering Board, European Coalition for AI in Fundamental Physics (2024–).

% ----------------------------------------------------------------------
\cvsection{Teaching and PhD Supervision and Examining}

\cvsmallheadingraw{Teaching at SISSA (2020–).}
Postgraduate courses: Bayesian Methods I \& II; Bayesian Inference and Machine Learning in Cosmology; Ethical Aspects of AI (Theoretical and Scientific Data Science PhD); AI and Information (Master in Science Communication).
\cvsmallheadingraw{Teaching at Imperial College London (2008–2020).}
Undergraduate courses: Physics of the Universe (3rd-year core);  Statistics of Measurement (2nd-year core); Cosmology (associate lecturer); Sun, stars and planets (associate lecturer).
\cvsmallheadingraw{Postgraduate schools}
Lecturer at 30+ international postgraduate schools and masterclasses (incl. Saas-Fee, ICTP, Carg\`ese, Heidelberg, Niels Bohr Institute, AIMS, ESAC); AI training course, SISSA, and EFC network AI training workshop, Bertinoro (2026).

\cvsmallheading{PhD supervision.}
11 completed PhD students (incl.\  two Springer Thesis Prize, two Winton Capital Best PhD Thesis prizes, RSS Emerging Applications Thesis Prize 2025); 2 ongoing PhDs; 14 postdocs (alumni now faculty at Cambridge, Uppsala, CONICET, Lawrence Livermore); 13 MSc/MSci students; 35 BSc projects; 19 external PhD/MSc examinations internationally.


% ----------------------------------------------------------------------
\cvsection{National and International Collaborations and Networks}

\cvsmallheading{Experimental collaborations.}
LSST Dark Energy Science Collaboration, member (2025–); XLZD Consortium for direct dark matter detection, member (2022–); DARWIN dark matter experiment, Executive Board member (2015–) and Science Working Package leader (2014–); ATLAS Collaboration, Short Term Associate (2014–16); XENON100, Associated Scientist (2014–18); Fermi-LAT, Affiliated Scientist (2012–15); Euclid ESA mission, Theory Working Group member (2010–14), likelihood Work Package co-lead (2020–22); co-author of the {\ttfamily SuperBayeS} public code for supersymmetric parameter analysis (2006–15).

\cvsmallheading{Networks and consortia.}
Centro Nazionale FAIR (PNRR), Spoke Edge and Exascale Computing (2022–26); Centro Nazionale HPC (PNRR), Spoke COSMOS (2022–26); Programme Committee Member, `ML and AI for Science', FVG Data Science Institute (2022–); active member, LSST/DESC mentoring programme (2026–); Management Committee Member, EU COST Action AIFORIA, Artificial Intelligence for Fundamental Physics (2026–).

\cvsmallheading{Advisory and coordination.}
External Advisory Board, McWilliams Center for Cosmology and Astrophysics, Carnegie Mellon University (2026–27); Scientific Consultant, ICTP Trieste (2022); Co-lead, Astrostatistics section of EuCAPT White Paper (2021); Astrostatistics section lead, Euclid Theory Working Group (2013, 2018); Scientific Advisor on Dark Energy probes, STFC UK (2005–07).

%\cvsmallheading{International partnerships.}
%Active collaborations across the UK (Cambridge, Imperial, Edinburgh), USA (KITP/UCSB, Penn State, Lawrence Livermore), Switzerland (Geneva, Bern, ETHZ), the Netherlands (Amsterdam), Spain (IFT Madrid, Valencia), South Africa (AIMS, Stellenbosch) and India (IUCAA Pune).


% ----------------------------------------------------------------------
\cvsection{Scientific Organisations and Coordination of Academic Activities}

\cvsmallheading{International executive and steering boards.}
International Executive Committee, INSAP conference series (2017–); Council of the International Astrostatistics Association (2015–); Outreach and Knowledge Transfer Leader, IAA (2015–); Astrostatistics Committee, International Statistical Institute (2011–); Working Group on Equity and Inclusion, Commission C1, IAU (2016–).

\cvsmallheading{Invited talks.}
30+ invited scientific talks and 110+ research talks and seminars, including plenary, EuCAIFCon, Heidelberg (2026); keynotes at Ital-IA (2025) and IFPU Focus Week (2025); colloquia at CERN (data science seminar, 2026), Ljubljana (2026), ICTP (2024), Z\"urich (2024), Bern (2017, 2024) and LMU Munich (2023).

\cvsmallheading{Conference and school organisation (selection from 45+ events).}
LOC Chair, Convegno Nazionale di Comunicazione della Scienza, Trieste (2025); SOC, EuCAIFCon (Amsterdam 2024; Cagliari 2025); SOC, INSAP XI–XIV (CalTech 2022, Corfu 2024, Belfast 2025, Rome 2026); SOC, Advanced Schools on Applied Machine Learning, ICTP (2024, 2026); SOC, Advanced School on Foundation Models for Scientific Discovery, ICTP (2025); Chair and founder, `Cosmostat' conference (Switzerland 2009; Banff 2013); Lead Organiser, KITP programme `Multi-probe Identification of Dark Matter', UCSB (2013).

\cvsmallheading{Evaluation and assessment panels.}
Evaluator: Emmanuel College Cambridge Research Fellowship (2025); Agence Nationale de la Recherche, France (2026); Swiss National Science Foundation (2024, 2025); ERC Consolidator Grants (2026); ERC Synergy Grants (2025); ERC Starting Grants (2021); UKRI Future Leaders (2025); Dutch Research Council (2024, 2026);  Canada Research Chairs programme (2020), South African National Research Foundation (2007–). 
Promotion and tenure assessor: Cambridge University  (2025); U.~Minnesota (2023); habilitation panel, U.~Nova Gorica (2026).
Panel member: STFC Rutherford Fellowships (2014–15); Royal Society International Grants (2007–12); Dutch Research Council ORC programme (2024, 2026).


% ----------------------------------------------------------------------
\cvsection{Editorial Activity}

\cvsmallheading{Editorial boards.}
Scientific Editor and founding member of the inaugural editorial board of {\itshape RAS Techniques and Instruments} (RASTI), open-access journal of the Royal Astronomical Society (2021–). Member, {\itshape Astronomy and Geophysics} Editorial Board (2009–12).

\cvsmallheading{Guest editor.}
{\itshape Statistical Analysis and Data Mining} (Wiley), special issue on Astrostatistics with A.~Gandy, vol.~6, 1–2 (2013); special issue, {\itshape Physics of the Dark Universe} (Elsevier, 2013).

\cvsmallheading{Peer review.}
80+ papers refereed for {\itshape Nature Physics}, {\itshape Physical Review Letters}, {\itshape Physical Review D/E}, {\itshape ApJ}, {\itshape MNRAS}, {\itshape JCAP}, {\itshape JHEP}, {\itshape Physics Letters B}, {\itshape New Astronomy}, {\itshape Computer Physics Communications}, {\itshape Science and Education}, {\itshape RASTI}.



% ----------------------------------------------------------------------
\cvsection{Membership of Scientific Societies}

\cvsmallheading{Elected fellow.}
European Coalition for AI in Fundamental Physics (EuCAIF, 2025); Higher Education Academy, UK (Senior fellow, SFHEA, 2021; Fellow, FHEA, 2011); International Astrostatistics Association (2016); Royal Astronomical Society (2005, FRAS).

\cvsmallheading{Academic}
Imperial College London Data Science Institute (2018–2022); Next Einstein Visiting Fellow, AIMS Cape Town (2011–2013).

\cvsmallheading{Professional member.}
International Society for Bayesian Analysis; International Astronomical Union; Swiss Society for Astronomy and Astrophysics; Royal Astronomical Society (UK); European Astronomical Society; International Statistical Institute, Astrostatistics Committee; Gresham Society; Royal Astronomical Society Dining Club.

\cvsmallheading{Affiliate.}
Chartered Management Institute, UK (2021–22).




% ----------------------------------------------------------------------
\cvsection{Other Work Experience and Consultancy}

\cvsmallheading{Advisory roles.}
Board of Directors, Fondazione Internazionale Trieste per il Progresso e la Libert\`a delle Scienze (2026–);  Scientific Advisory Board, RACHAEL S.r.l.\ (2021–; Chairman of the Board 2021–24).

\cvsmallheading{Entrepreneurship.}
Co-founder and Director, {\itshape Data Fusion Consultants Ltd} (2012–20); Non-executive Director, Imperial Consultants (2018–20).

\cvsmallheading{Industry and media clients (selection).}
Beluga Group, Savio Macchine Tessili, Oxford University Press, Nature Publishing Group, Red Planet Pictures (BBC `Death in Paradise'), Left Bank Pictures (Netflix `Origin'), Fantastic Productions (Netflix), BBC Sky at Night, London Science Museum, Deutsches Museum Munich, Museo del Castello di Miramare, Regione Friuli Venezia Giulia, Institute of Physics. Statistical consultancy for clients in shipping, energy, property, banking, fashion and online commerce.

\cvsmallheading{Executive training.}
Imperial College Business School Executive Education (Sberbank Digital Technology, 2020–21; Naspers Executives Programme, 2019); Generali Business Translators (2022); CINECA Academy (2022); Intesa San Paolo (2024); FINECO Advisory Seminar keynote (2025); Generali Security Council keynote (2026); Machine Learning for Business Analytics, MBA, MIB Trieste (2023–).

% ----------------------------------------------------------------------
\cvsection{Public Engagement and Science Communication}

\cvsmallheading{Books for the general public.}
{\itshape Starborn: How the Stars Made Us} (Basic Books, 2023): BBC Radio 4 Book of the Week, BBC Sky at Night 5-star review, {\itshape New Scientist} `22 Best Non-Fiction and Popular Science Books of 2023', {\itshape Smithsonian Magazine} `Ten Best Science Books of 2023', 2025 Nautilus Book Awards Gold Winner; translated into Italian, Korean, Spanish and Chinese. {\itshape The Edge of the Sky} (Basic Books, 2014): cosmology in the 1{,}000 most common English words; named {\itshape Foreign Policy} `Global Thinker 2014';  `Great Popular Science Books 2014', {\em The Huffington Post}; `Favorite Physics Books of 2014', {\em Scientific American};
translated into Korean, German and Catalan.


\cvsmallheading{Public reach.}
\refcvcount{nPublicLectures}+ public lectures and talks at major science festivals (Cheltenham, Edinburgh, New Scientist Live, Royal Albert Hall, Royal Institution, Mantova, Torino, Genova); \refcvcount{nAllPEPubs}+ articles for the general public ({\itshape TIME}, {\itshape The Guardian}, {\itshape New Scientist}, {\itshape La Stampa}, {\itshape Le Scienze}); broadcast appearances on BBC Radio 4, BBC World Service, NPR, RAI, RSI Swiss National Radio/TV, Sky TG24.

\cvsmallheading{Art and science projects.}
Scientific advisor and project lead, {\itshape La storia dell'Universo}, permanent exhibit, SISSA (2026); scientific advisor, {\itshape The unwetlands}, satellite-based investigation of Italy's lost wetlands (2026); co-author and creative director, {\itshape LIBRA} (2021, multimedia theatre play, performed at Miramare Castle, Reggio Calabria, Gorizia 2025 European Capital of Culture); AI consultant, {\itshape Parl-IA-moci} (2023, theatre play with Diana H\"obel and ChatGPT); {\itshape A Multi-Sensorial Dark Matter Experience} (2019); {\itshape Dinner with Copernicus} (2018); {\itshape g-ASTRONOMY} multi-sensorial cosmology dinners; {\itshape The Theory of Everything} short film (audience award, International Short Film Festival Berlin 2013).

\cvsmallheading{Consulting for film and media.}
Scientific consultant for Red Planet Pictures (BBC `Death in Paradise'), Left Bank Pictures (`Origin', Netflix), Fantastic Productions (Netflix feature film), Nature Publishing Group, Oxford University Press.

\cvsmallheading{Jury and awards panels.}
Jury Member, Premio Cosmos (Italian national science book award, 2022–); International FameLab semi-final judge, Cheltenham (2015).
